\section{Carico e distribuzione del lavoro}
\label{caricoLavoro}
Come precedentemente anticipato, ogni membro ha contribuito in qualche modo al completamento di tutti i \textit{deliverable}. Di seguito viene descritto il carico di lavoro in numero di ore per ciascun membro del gruppo suddiviso per \textit{deliverable}.

\begin{table}[H]
    \centering
    \begin{tabular}{|c|c|c|c|c|c|}
        \hline
        & \texttt{D1} & \texttt{D2} & \texttt{D3} & \texttt{D4} & \textbf{Totale}\\
        \hline
        \textbf{Alessio Pesaresi} & 35 & 13 & 36 & 8 & 92\\
        \hline
        \textbf{Michele Porcellato} & 35 & 80 & 15 & 6 & 136\\
        \hline
        \textbf{Marco Piovesan} & 66 & 17 & 37 & 10 & 130\\
        \hline
        \textbf{Totale} & 136 & 110 & 88 & 24 & 358\\
        \hline
    \end{tabular}
\end{table}

Alcuni squilibri che si possono notare soprattutto nei \textit{deliverable} \texttt{D2} sono dovuti al fatto che le più avanzate competenze precedentemente acquisite da \textit{Michele Porcellato} in materia di sviluppo web e di \textit{Marco Piovesan} in materia di progettazione hanno permesso di completare più velocemente la parte di sviluppo e progettazione rispettivamente. Il carico di lavoro di \textit{Alessio Pesaresi} è stato più equilibrato rispetto agli altri due membri del gruppo in quanto la sua attività orizzontale ha richiesto un impegno costante in tutti i \textit{deliverable}.